% Example: Inline math ($...$) vs display math (\[...\])
% Demonstrates the difference between math embedded in text and math
% displayed on its own centered line.
\documentclass[12pt, a4paper]{article}
\usepackage{amsmath, amssymb, amsthm}
\usepackage{fontspec}
\begin{document}

\section*{Inline vs Display Math}

Inline math flows within the surrounding text. For example, the famous
identity \( e^{i\pi} + 1 = 0 \) relates five fundamental constants, and
the quadratic formula \( x = \frac{-b \pm \sqrt{b^2 - 4ac}}{2a} \) can
appear mid-sentence like this.

Display math is typeset on its own centered line, which improves
readability for larger expressions:

\[
  \int_{-\infty}^{\infty} e^{-x^2}\,dx = \sqrt{\pi}
\]

Another display example:

\[
  \sum_{n=1}^{\infty} \frac{1}{n^2} = \frac{\pi^2}{6}
\]

Notice how inline math like \( a^2 + b^2 = c^2 \) stays compact, while
the same equation in display mode stands out:

\[
  a^2 + b^2 = c^2
\]

\end{document}
